\section{Protection Status: Establishing the Absence of Commercial Anti-Tamper}
\label{sec:denuvo}

This section answers the question the reader most likely arrived for: \emph{how
do you determine, from a binary alone, that a commercial anti-tamper system is
not present?} The answer has a strong form and a weak form, and most published
methodology uses only the weak form.

\subsection{Why the question had to be answered first}

Commercial anti-tamper systems of the Denuvo class operate by taking selected
regions of the program's control flow and replacing them with an interpreter for
a private instruction set, generated per build. The practical consequences for
this project were three:

\begin{enumerate}
\item \textbf{Compatibility.} Such a layer interacts closely with operating
      system internals. A compatibility layer that reimplements those internals
      may or may not satisfy it, and the failure mode is not usually a clean
      error message.
\item \textbf{Effort allocation.} If the protection were present, the ceiling on
      this project might be architectural rather than engineering---in which case
      the correct decision is to stop, as was in fact taken for the contrasting
      title discussed in Section~\ref{sec:contrast}.
\item \textbf{Threat model.} The outstanding malware theory from
      Section~\ref{sec:malware} was that the protection had been
      \emph{reimplemented} by a third party, and that this reimplementation was
      the injection vector. Proving no such layer exists retires that theory
      entirely.
\end{enumerate}

The subject of analysis is \verb|Retail\007FirstLight.exe|, 64{,}228{,}744 bytes,
MD5 \verb|3e0c3e10a9fdb15397f5bc50aeb5e09b|, extracted from the fifth container.

\subsection{The strong argument: signature validation}
\label{sec:sig}

Windows executables may carry an Authenticode signature~\cite{authenticode}: a
PKCS\#7 / CMS structure~\cite{rfc5652} embedded in the certificate directory of
the PE file~\cite{pecoff}, containing a cryptographic digest computed over the
file's contents with three regions excluded (the checksum field, the certificate
table entry, and the certificate data itself).

Validation proceeds in two independent steps, and the distinction between them is
the crux of this section:

\begin{enumerate}
\item \textbf{Chain validation} asks whether the signing certificate descends
      from a trusted root. This says who signed it.
\item \textbf{Digest validation} recomputes the Authenticode hash over the
      \emph{file as it exists now} and compares it to the digest stored inside
      the signature blob. This says whether the file has changed since it was
      signed.
\end{enumerate}

For this executable, both hold. The recomputed SHA-256 equals the digest embedded
in the PKCS\#7 structure. The signer is \textbf{IO~INTERACTIVE~A/S}, via
GlobalSign GCC R45 CodeSigning CA 2020, with a signing time of
2026-07-20 15:43:44 UTC.

\subsubsection{Why this settles the question}

The reasoning is short enough to state completely:

\begin{quote}
The signature validates over the current bytes. Therefore the current bytes are
the bytes the publisher signed. Therefore \textbf{nobody has modified this
executable}---not to add anything, and, crucially, \textbf{not to remove
anything either}. Anti-tamper protection of the class in question is applied to
the executable at build time by the publisher and lives inside the executable's
own sections. If it had ever been present in this build, it would still be
present now, and it would be visible. It is not visible. Therefore it was never
applied to this build.
\end{quote}

This is a \emph{sufficient} argument, and it is the strong form. It requires no
dynamic analysis, no debugger, no emulator, and no unpacking. It is also
self-verifying in a way that heuristics are not: a forged conclusion would
require a hash collision against SHA-256.

It carries a second payload for free. The unresolved malware theory---that the
protection was reimplemented and used as an injection vector---requires
modification of the executable. The signature forecloses that. \textbf{One
measurement answers both questions.}

\subsubsection{The scope of the claim, stated honestly}

The argument establishes that \emph{this executable} is unmodified and
unprotected. It does not establish anything about the other 10{,}949 files in the
package; those were covered separately, by the differential comparison in
Section~\ref{sec:malware}. It also does not establish that no protection exists
anywhere in the product---a separate protected loader could in principle exist as
a distinct binary. Section~\ref{sec:corrob} addresses that by enumerating what
the executable actually imports and what else ships beside it.

\subsection{Corroboration: three independent weak arguments}
\label{sec:corrob}

The strong argument is sufficient, but a single-instrument conclusion is a fragile
conclusion. Three further lines were run, each of which would independently have
raised suspicion had the protection been present.

\subsubsection{Entropy profiling}

Compressed, encrypted, and virtualised code is high-entropy; ordinary compiled
code is not. Entropy analysis is the standard first-pass instrument for detecting
packing~\cite{lyda2007}, with the well-documented caveat that it produces both
false positives and false negatives in the general case~\cite{sokpacker}.

Rather than compute a single figure for the file, which averages away exactly the
signal of interest, the \verb|.text| section was profiled in 43 windows of 1\,MiB
each.

\begin{table}[htbp]
\caption{Shannon entropy profile of the game executable}
\label{tab:entropy}
\centering
\footnotesize
\begin{tabularx}{\columnwidth}{@{}L{3.3cm}X@{}}
\toprule
\textbf{Measurement} & \textbf{Value} \\
\midrule
Whole file & 6.354 \\
\verb|.text|, 43 windows of 1\,MiB & range 5.906--6.700, mean 6.505 \\
Windows at or above 7.5 & \textbf{0} \\
Expected for a virtualised region & 7.9 and above \\
\bottomrule
\end{tabularx}
\end{table}

The distribution is tight, unimodal, and entirely consistent with optimised
native x86-64. There is no high-entropy island anywhere in the code section. Had
a virtualising layer been present, it would appear as a contiguous run of windows
well above 7.5, and windowing is precisely what makes it visible: a whole-file
figure of 6.354 could in principle conceal a small high-entropy region, but 43
individually reported windows cannot.

\subsubsection{Section and import table analysis}

\begin{table}[htbp]
\caption{Structural indicators}
\label{tab:pe}
\centering
\footnotesize
\begin{tabularx}{\columnwidth}{@{}L{3.0cm}X@{}}
\toprule
\textbf{Indicator} & \textbf{Observation} \\
\midrule
Sections & \verb|.text|, \verb|.rdata|, \verb|.data|, \verb|.pdata|,
            \verb|_RDATA|, \verb|.rsrc|, \verb|.reloc| --- stock MSVC layout \\
Anomalous sections & none \\
Imports & 52 DLLs, 753 symbols; a normal compiler-emitted table \\
Runtime & ucrt and \verb|VCRUNTIME140| \\
Protection strings & zero hits for the vendor name, VMProtect, or Themida \\
Size profile & the several-hundred-megabyte bloat characteristic of a protected
               build is absent \\
\bottomrule
\end{tabularx}
\end{table}

Protected builds are structurally distinctive. They typically carry an extra
section with a non-standard name, a truncated or stub import table with the real
imports resolved at runtime, and a substantial size increase over the unprotected
build. None of these is present. The import table in particular is decisive in
its own small way: a protected binary usually cannot afford a full static import
table, because the protection needs to mediate library resolution.

\subsubsection{What the DRM actually is}

Having established what is absent, the analysis identified what is present. The
only protection mechanism in the package is Steamworks licensing, and the
countermeasure is a third-party Steam API emulator.

\begin{table}[htbp]
\caption{The actual protection and countermeasure}
\label{tab:rune}
\centering
\footnotesize
\begin{tabularx}{\columnwidth}{@{}L{3.0cm}X@{}}
\toprule
\textbf{File} & \textbf{Identification} \\
\midrule
\verb|steam_api64.rne| & 319{,}584 B --- the genuine publisher-shipped Valve
                         library, renamed \\
\verb|steam_api64.dll| & 1{,}098{,}392 B --- the third-party replacement \\
\verb|steam_emu.ini| & configuration; carries the author's banner, the
                       application ID, and a configuration section \\
\bottomrule
\end{tabularx}
\end{table}

The mechanism, characterised but not extended (Section~\ref{sec:ethics}), is a
library-loading hook that looks for renamed twins of each library. The
architecturally relevant point is that \textbf{it replaces only Valve's own API
surface}---a documented interface between the game and a storefront---and does
not touch the game's code. This is consistent with, and independently confirms,
the signature finding.

\subsection{The contrasting case, and why it terminated}
\label{sec:contrast}

A second title on the same machine provides the control that makes the above
meaningful: it \emph{does} carry the commercial anti-tamper system, still present
and active. Its accompanying countermeasure is not a library replacement but a
\textbf{hypervisor}: a pair of kernel drivers implementing AMD-V and Intel VT-x
support respectively, which load beneath the Windows kernel and answer the
protection's queries from a layer below it. Source for the virtualisation
component ships in the archive.

This approach cannot work on Linux under a compatibility layer, for two
independent reasons, either of which is sufficient:

\begin{enumerate}
\item \textbf{There is no kernel to load into.} Wine implements Windows APIs in
      user space; it does not implement a Windows kernel, so there is no object
      into which a \verb|.sys| driver can be loaded.
\item \textbf{The hardware is already taken.} The CPU's virtualisation extensions
      are owned by the Linux host. A second hypervisor cannot claim them.
\end{enumerate}

We note this explicitly as an \emph{architectural impossibility}, in the narrow
sense argued in Section~\ref{sec:lessons}: not ``this would be difficult'' but
``this specific mechanism has no counterpart on this platform.'' The
investigation of that title stopped there, which is the correct outcome and the
whole value of asking the protection question before the engineering question.

\subsection{Summary of the method}

For a reader wishing to reproduce this determination on another binary, the
procedure in decreasing order of decisiveness:

\begin{enumerate}
\item \textbf{Validate the Authenticode digest over current bytes.} If it
      validates, the file is the publisher's build and the question is answered.
      Do not skip to entropy; this step is both cheaper and stronger.
\item \textbf{Profile entropy in windows}, not whole-file. Report the window
      count above threshold, and state the threshold.
\item \textbf{Enumerate sections and imports.} A full, compiler-shaped import
      table is difficult for a protection layer to preserve.
\item \textbf{Identify what protection \emph{is} present.} ``No Denuvo'' is not
      ``no DRM,'' and conflating them produces a wrong prediction about runtime
      behaviour.
\end{enumerate}

For this title the answer is unambiguous: \textbf{there is no commercial
anti-tamper system, and there never was one in this build.} Nothing was
stripped, because there was nothing to strip. The engineering path was therefore
open, and the remainder of this paper is about walking it.
